{\ttfamily
    1. ¿En qué contextos ha desempeñado actividades como formación, planificación, ejecución o supervisión de mantenimiento preventivo?\\

    2. ¿Cuáles considera que son los conceptos o conocimientos de mantenimiento preventivo más relevantes que ha aplicado o presenciado?\\

    3. ¿Cuál es su experiencia con el uso de tecnologías de la industria 4.0 y en qué ámbito del mantenimiento ha presenciado la aplicación de las mismas?\\

    4. ¿Conoce o es partícipe de algún caso de uso de la tecnología de realidad aumentada en el mantenimiento?\\

    5. ¿Qué situaciones problemáticas ha experimentado o identificado que se relacionan con el mantenimiento preventivo que pueden mejorarse?\\

    6. ¿Qué aspectos del mantenimiento preventivo considera que se pueden mejorar con la implementación de alguna tecnología de la información y comunicación?\\

    7. ¿Qué ventajas percibe en permitir que los operarios de mantenimiento gestionen la información en un dispositivo móvil como tablet o smartphone?\\

    8. ¿Cómo afectaría la inclusión de dispositivos móviles al flujo de trabajo de los operarios en un proceso de mantenimiento preventivo?\\

    9. ¿Qué tipo de información le resultaría más relevante tener en tiempo real durante una intervención de mantenimiento?\\

    10. ¿Cómo considera que podría ayudar una herramienta basada en realidad aumentada en la reducción de errores de ejecución de los distintos procesos de mantenimiento?\\

    11. ¿Qué funcionalidades considera que traerán beneficios en el uso de la realidad aumentada a la asistencia del mantenimiento preventivo y cuáles serían sus beneficios?\\

    12. ¿Qué aspectos considera usted que se mejoran con la aplicación de la realidad aumentada en el mantenimiento industrial?\\

    13. ¿Cuáles son los mayores retos que se anticipan en la capacitación de personal para el uso de herramientas de realidad aumentada y de gestión de la información en dispositivos móviles?\\

    14. ¿Considera que el uso de la realidad aumentada puede ayudar en la capacitación del personal en tareas de mantenimiento, en favor de disminuir los tiempos de dichas capacitaciones y mejorar el tiempo de asimilación de dichas rutinas o procesos?\\
}